% настройки polyglossia
\setdefaultlanguage{russian}
\setotherlanguage{english}

% локализация
\addto\captionsrussian{%
  \renewcommand{\figurename}{Рисунок}%
  \renewcommand{\partname}{Глава}%
  \renewcommand{\contentsname}{\centerline{Содержание}}%
  \renewcommand{\listingscaption}{Листинг}%
}

% основной шрифт документа (Times New Roman)
% \setmainfont{CMU Serif}
% \newfontfamily\cyrillicfont{CMU Serif}[Script=Cyrillic]
\setmainfont[Ligatures={TeX,Historic}]{Times New Roman}
\defaultfontfeatures{Ligatures={TeX},Renderer=Basic}

% URL тем же шрифтом, что и основной текст
\renewcommand{\UrlFont}{\normalfont}

% аккуратные переносы внутри URL
\def\UrlBreaks{%
  \do\/\do-\do_\do.\do?\do&\do=\do:\do\%\do\#%
}

% оформление полей библиографии (biblatex-gost)
% всё — обычным шрифтом Times New Roman
\DeclareFieldFormat*{doi}{\textnormal{#1}}
\DeclareFieldFormat*{eprint}{\textnormal{#1}}
\DeclareFieldFormat*{title}{\textnormal{#1}}
\DeclareFieldFormat*{journaltitle}{\textnormal{#1}}
\DeclareFieldFormat*{booktitle}{\textnormal{#1}}
\DeclareFieldFormat*{note}{\textnormal{#1}}

% URL — через \url, чтобы работали \UrlBreaks
\DeclareFieldFormat{url}{\url{#1}}

% перечень использованных источников
\addbibresource{refs.bib}

% настройка полей страницы
\geometry{top=2cm}
\geometry{bottom=2cm}
\geometry{left=3cm}
\geometry{right=15mm}
\geometry{bindingoffset=0cm}

% настройка ссылок и метаданных документа
\hypersetup{%
  unicode=true,
  colorlinks=true,
  breaklinks=true,
  linkcolor=black,
  citecolor=black,
  filecolor=black,
  urlcolor=black
}

% настройка подсветки кода и окружения для листингов
\usemintedstyle{colorful}
\newenvironment{code}{\captionsetup{type=listing}}{}

% шрифт для листингов с лигатурами
\setmonofont{FiraCode-Regular.otf}[%
  SizeFeatures={Size=10},%
  Path=templates/,%
  Contextuals=Alternate%
]

% подписи рисунков
\captionsetup[figure]{%
  labelsep=space,%
  justification=centering,%
  singlelinecheck=false%
}

% подписи таблиц
\DeclareCaptionFormat{hfillstart}{\hfill#1#2#3\par}
% \captionsetup[table]{format=hfillstart,labelsep=newline,justification=centering,skip=-10pt,textfont=bf}
\captionsetup[table]{%
  name={Таблица},%
  labelsep=space,%
  justification=raggedright,%
  singlelinecheck=false,%
  font=normalfont,%
  labelfont=normalfont,%
  skip=6pt%
}

% путь к каталогу с рисунками
\graphicspath{{fig/}}

% учёт титульного листа
\makeatletter
\renewenvironment{titlepage}{%
  \thispagestyle{empty}%
}{}
\makeatother

% нумерация рисунков и таблиц по разделам
\counterwithin{figure}{section}
\counterwithin{table}{section}

% формат заголовков \title
\titlelabel{\thetitle.\quad}

% математика
\mathtoolsset{showonlyrefs=true}
\theoremstyle{plain}
\newtheorem{theorem}{Теорема}[section]
\newtheorem{proposition}[theorem]{Утверждение}
\theoremstyle{definition}
\newtheorem{corollary}{Следствие}[theorem]
\newtheorem{problem}{Задача}[section]
\theoremstyle{remark}
\newtheorem*{nonum}{Решение}

% настоящее матожидание
\newcommand{\MExpect}{\mathsf{M}}

% объявили оператор!
\DeclareMathOperator{\sgn}{\mathop{sgn}}

% перенос знаков в формулах (по Львовскому)
\newcommand*{\hm}[1]{#1\nobreak\discretionary{}{\hbox{$\mathsurround=0pt #1$}}{}}

% настройки эпиграфов
\setlength{\epigraphwidth}{0.5\textwidth}
\setlength{\epigraphrule}{0pt}


